\section{Related Work}
\label{sec:related}

MAgHARCM sits at the intersection of LLM-based code translation,
program-repair agents, and neuro-symbolic migration tooling. The four
closest systems are ReCodeAgent, AlphaTrans, CodePlan, and Oxidizer;
MatchFixAgent and TRAM round out the corpus because they share the
same agent-graph or repair-loop structure.

\paragraph{ReCodeAgent.}
ReCodeAgent~\citep{ibrahimzada2025recodeagent} is the architectural
template MAgHARCM extends: a four-agent Analyzer / Planning /
Translator / Validator cyclic graph over a Model-Context-Protocol tool
layer. We adopt the same decomposition verbatim and target the same
public benchmarks (GildedRose, the Oxidizer corpus, the AlphaTrans
Java-Python subject set). MAgHARCM differs from ReCodeAgent in three
concrete directions: it is wired against a local Ollama runtime rather
than a hosted LLM, it provides a tree-sitter-native LSP provider with
graceful ABCoder-MCP fallback rather than requiring ABCoder, and it
adds a coverage-guided characterization-test pass on top of the
existing Validator. The Sample 1 results in Section~\ref{sec:eval} are
the first end-to-end exercise of the four-agent graph on a 4B coding
model.

\paragraph{AlphaTrans.}
AlphaTrans~\citep{alphatrans2024} is the Java-Python translation
baseline we used to seed Sample 3. AlphaTrans frames translation as a
structured type-and-API mapping problem and applies a sequence of
neuro-symbolic passes to reconcile Java generics with Python's dynamic
type system. MAgHARCM does not yet have a Java-Python sample run; the
Sample 3 slot in Table~\ref{tab:samples} is reserved for it.

\paragraph{CodePlan.}
CodePlan is a plan-then-translate neuro-symbolic baseline that emits
an explicit dependency-ordered execution plan before invoking the
LLM. ReCodeAgent's Planning agent borrows this idea; MAgHARCM's
Planning agent inherits it. We treat CodePlan as a second-order
reference and use the Oxidizer/stats corpus (originally curated for
CodePlan and Oxidizer) as the Sample 2 subject.

\paragraph{Oxidizer.}
Oxidizer~\citep{oxidizer2024} is the Go-Rust translation baseline.
Sample 2 (\texttt{stats}, a 73-file Go statistics library) and
Sample 4 (\texttt{gohistogram}, a smaller Go histogram library) are
both drawn from the Oxidizer subject set. The 517-fragment ceiling we
observe on Sample 2 is, in retrospect, exactly the failure mode
Oxidizer's static-analysis-first design avoids: by pre-computing a
typed cross-language mapping it never asks the LLM to ingest a
517-fragment single prompt. This is the same problem NEW-PRIM-23
addresses from the agent side.

\paragraph{MatchFixAgent.}
MatchFixAgent~\citep{matchfixagent2024} is a repair-agent baseline that fixes
compiler diagnostics one error at a time. We cite it because the
Translator's repair mode in MAgHARCM is a direct descendant of the
MatchFixAgent loop: a Validator diagnostic is fed back to the
Translator as a structured prompt, and the Translator emits a
single-file patch. The MAgHARCM loop differs in that it is
multi-file (the Planner can ask the Translator to edit any number of
files in one iteration) and that the Translator emits a full file
rewrite rather than a unified-diff patch.

\paragraph{TRAM.}
TRAM~\citep{tram2024} is a multi-agent translation-and-repair baseline
that separates translation from test generation, the same dichotomy
that ReCodeAgent's Validator formalises. The TRAM corpus and the
ReCodeAgent corpus overlap heavily; MAgHARCM's characterization-test
pass is most directly traceable to TRAM's test-generation stage.

\paragraph{Local-SLM agent harnesses.}
MAgHARCM is not the first system to drive a multi-agent code task
against a local SLM. We treat this as an engineering decision (no
API access on the target host, predictable cost) rather than a
research contribution, and we cite the local-model agent frameworks
that informed the harness layout without claiming novelty over them.
